\chapter{Bilangan Besar}\label{ch:g4:numbers}

Ribuan adalah batas atas pada \cref{ch:g3:numbers}; sekarang batas
itu naik sampai \emph{jutaan}. Trik yang membuat bilangan raksasa
mudah dibaca adalah memotong angkanya menjadi kelompok tiga --- yaitu
\omterm{def:g4:numbers:classes}{kelas}.

\section{Kelas tiga angka}

\begin{definition}[Kelas]\label{def:g4:numbers:classes}
Mulai dari kanan, angka sebuah bilangan dikelompokkan menjadi
\emph{kelas} yang terdiri atas tiga angka: kelas \emph{satuan}, kelas
\emph{ribuan}, dan kelas \emph{jutaan}. Setiap kelas punya satuan,
puluhan dan ratusannya sendiri. Misalnya
\[
6\,309\,452 = 6 \text{ jutaan} + 309 \text{ ribuan}
+ 452 \text{ satuan},
\]
dibaca ``enam juta tiga ratus sembilan ribu
empat ratus lima puluh dua''.
\end{definition}

\begin{omfigure}
\begin{tikzpicture}[scale=0.86]
  \draw (0,0) grid[xstep=1.15, ystep=0.75] (10.35,1.5);
  \node[font=\footnotesize, omThm] at (1.725,1.85) {jutaan};
  \node[font=\footnotesize, omDef] at (5.175,1.85) {ribuan};
  \node[font=\footnotesize, omMeth] at (8.625,1.85) {satuan};
  \foreach \x/\l in {0.575/r, 1.725/p, 2.875/s,
                     4.025/r, 5.175/p, 6.325/s,
                     7.475/r, 8.625/p, 9.775/s}
    \node[font=\footnotesize] at (\x,1.12) {\l};
  \foreach \x/\d in {2.875/6, 4.025/3, 5.175/0, 6.325/9,
                     7.475/4, 8.625/5, 9.775/2}
    \node[font=\large] at (\x,0.37) {$\d$};
  \draw[omThm, thick] (0,0) rectangle (3.45,1.5);
  \draw[omDef, thick] (3.45,0) rectangle (6.9,1.5);
  \draw[omMeth, thick] (6.9,0) rectangle (10.35,1.5);
\end{tikzpicture}

{\small Tabel \omterm{def:g4:numbers:classes}{kelas} dari $6\,309\,452$ (r $=$ ratusan, p $=$
puluhan, s $=$ satuan pada tiap \omterm{def:g4:numbers:classes}{kelas}). Membaca bilangan besar
berarti membaca setiap \omterm{def:g4:numbers:classes}{kelas}, lalu menyebut namanya.}
\end{omfigure}

\begin{example}[Menulis dengan kelas]\label{ex:g4:numbers:write}
``Dua belas juta empat puluh ribu lima ratus'' punya $12$ di \omterm{def:g4:numbers:classes}{kelas}
jutaan, $040$ di \omterm{def:g4:numbers:classes}{kelas} ribuan, dan $500$ di \omterm{def:g4:numbers:classes}{kelas} satuan:
\[
12\,040\,500 .
\]
Angka \omterm{def:g1:counting:zero}{nol} pada $040$ itu penting: setiap \omterm{def:g4:numbers:classes}{kelas} (kecuali yang paling
kiri) harus menampilkan ketiga angkanya.
\end{example}

\begin{example}[Berapa ribuan?]\label{ex:g4:numbers:contained}
Bilangan $6\,309\,452$ \emph{memuat} $6\,309$ ribuan (semua yang ada
di kiri \omterm{def:g4:numbers:classes}{kelas} satuan): $6\,309\,452 = 6\,309 \times
1\,000 + 452$. Bertanya ``berapa ribuan'' bukan menanyakan satu
angka, melainkan seluruh awalan bilangan itu.
\end{example}

\section{Membandingkan dan membulatkan}

\begin{method}[Membandingkan bilangan besar]\label{met:g4:numbers:compare}
Persis seperti pada \cref{met:g3:numbers:compare}, dengan \omterm{def:g4:numbers:classes}{kelas}
sebagai jalan pintas: angka yang lebih banyak menang; kalau
panjangnya sama, bandingkan dari kiri --- \omterm{def:g4:numbers:classes}{kelas} demi \omterm{def:g4:numbers:classes}{kelas}, lalu
angka demi angka.
\end{method}

\begin{example}\label{ex:g4:numbers:compare}
$2\,145\,890$ dan $2\,148\,540$: jutaannya sama ($2$), \omterm{def:g4:numbers:classes}{kelas}
ribuannya $145 < 148$, jadi
\[
2\,145\,890 < 2\,148\,540 .
\]
\end{example}

\begin{definition}[Pembulatan]\label{def:g4:numbers:round}
\emph{Membulatkan}\index{pembulatan} sebuah bilangan ke ribuan
terdekat (atau ke puluhan, ratusan, jutaan\,\dots) berarti
menggantinya dengan kelipatan $1\,000$ yang paling dekat. Jika
bilangannya berada tepat di tengah, bulatkan ke atas. Dalam praktik:
lihatlah angka tepat di bawah tempat pembulatan --- $5$ atau lebih
dibulatkan ke atas, $4$ atau kurang dibulatkan ke bawah.
\end{definition}

\begin{example}\label{ex:g4:numbers:round}
$37\,614$ \omterm{def:g4:numbers:round}{dibulatkan} ke ribuan terdekat: angka ratusannya $6$, jadi
\omterm{def:g4:numbers:round}{bulatkan} ke atas: $38\,000$. \omterm{def:g4:numbers:round}{Dibulatkan} ke ratusan terdekat:
$37\,600$. \omterm{def:g4:numbers:round}{Dibulatkan} ke puluhan ribu terdekat: $40\,000$.
\end{example}

\begin{omfigure}
\begin{tikzpicture}[scale=1.05]
  \draw[-Stealth] (-0.3,0) -- (10.8,0);
  \draw (0,-0.1) -- (0,0.1);
  \node[below, font=\small] at (0,-0.12) {$37\,000$};
  \draw (10,-0.1) -- (10,0.1);
  \node[below, font=\small] at (10,-0.12) {$38\,000$};
  \draw (5,-0.08) -- (5,0.08);
  \node[below, font=\footnotesize] at (5,-0.12) {$37\,500$};
  \fill[omThm] (6.14,0) circle (2pt)
    node[above, font=\small] {$37\,614$};
  \draw[-Stealth, omDef, thick] (6.14,0.42) .. controls (8,0.85) ..
    (9.9,0.28);
\end{tikzpicture}

{\small \omterm{def:g4:numbers:round}{Membulatkan} $37\,614$ ke ribuan terdekat: letaknya sudah
melewati tanda tengah $37\,500$, jadi ia \omterm{def:g4:numbers:round}{dibulatkan} naik ke $38\,000$.}
\end{omfigure}

\begin{example}[Orde besaran di dunia nyata]\label{ex:g4:numbers:world}
\omterm{def:g4:numbers:round}{Pembulatan} adalah cara bilangan besar benar-benar dipakai: sebuah
negara berpenduduk $67\,842\,591$ jiwa disebut ``sekitar $68$ juta'';
sebuah kota berpenduduk $214\,376$ jiwa disebut ``sekitar
$200\,000$''. Angka persisnya toh berubah setiap hari --- angka yang
dibulatkanlah yang membawa keterangan yang berguna.
\end{example}

\section{Latihan}

\begin{exercise}[$\star$]\label{exo:g4:numbers:1}
Tulislah dengan angka, dengan \omterm{def:g4:numbers:classes}{kelas} yang dipisahkan: ``tiga juta
lima ratus dua puluh ribu''; ``tujuh juta delapan''; ``dua ratus
empat puluh ribu enam puluh''.
\end{exercise}

\begin{exercise}[$\star$]\label{exo:g4:numbers:2}
Tulislah dengan kata: $5\,204\,300$; \quad $10\,050\,000$; \quad
$999\,999$.
\end{exercise}

\begin{exercise}[$\star$]\label{exo:g4:numbers:3}
Pada $8\,472\,156$: angka mana yang menjadi puluhan ribu? Ratusan
pada \omterm{def:g4:numbers:classes}{kelas} jutaan? Berapa \emph{ribuan} yang dimuat bilangan itu
(lihat \cref{ex:g4:numbers:contained})?
\end{exercise}

\begin{exercise}[$\star$]\label{exo:g4:numbers:4}
Salin dan lengkapi dengan $<$ atau $>$:
\[
980\,000 \;?\; 1\,020\,000, \qquad
3\,456\,000 \;?\; 3\,465\,000, \qquad
7\,090\,050 \;?\; 7\,090\,500 .
\]
\end{exercise}

\begin{exercise}[$\star$]\label{exo:g4:numbers:5}
Urutkan dari yang terkecil ke yang terbesar: $2\,300\,000$; \ $980\,500$; \
$2\,030\,000$; \ $2\,299\,999$.
\end{exercise}

\begin{exercise}[$\star$]\label{exo:g4:numbers:6}
Bulatkan $63\,782$: ke puluhan terdekat; ke ratusan terdekat; ke
ribuan terdekat; ke puluhan ribu terdekat.
\end{exercise}

\begin{exercise}[$\star$]\label{exo:g4:numbers:7}
Bulatkan ke jutaan terdekat: $4\,721\,000$; \quad $6\,499\,999$;
\quad $9\,500\,000$.
\end{exercise}

\begin{exercise}[$\star$]\label{exo:g4:numbers:8}
Bilanglah dengan langkah: tiga langkah $1\,000$ dari $998\,500$; dua
langkah $100\,000$ dari $950\,000$. Apa yang terjadi pada kelasnya
ketika kamu menyeberangi satu juta?
\end{exercise}

\begin{exercise}[$\star$]\label{exo:g4:numbers:9}
Sebuah stadion memuat $81\,338$ penonton. Sebuah koran menulis
``lebih dari $80\,000$ penggemar''; sebuah radio berkata ``sekitar
$81\,000$''. Apakah keduanya benar? Masing-masing \omterm{def:g4:numbers:round}{membulatkan} ke tempat yang mana?
\end{exercise}

\begin{exercise}[$\star\star$]\label{exo:g4:numbers:10}
Dengan memakai masing-masing angka $1$ sampai $7$ tepat sekali,
tulislah bilangan tujuh angka yang terbesar, yang terkecil, dan yang
paling dekat dengan $4\,000\,000$.
\end{exercise}

\begin{exercise}[$\star\star$]\label{exo:g4:numbers:11}
Berapa bilangan bulat terbesar yang menjadi $50\,000$ ketika
\omterm{def:g4:numbers:round}{dibulatkan} ke ribuan terdekat? Dan berapa yang terkecil?

\end{exercise}
